\centerline{\bf \Huge Комбинаторный разнобой 1}

\problem{1}
В ряд выложены апельсины, бананы, яблоки и груши так, что все четыре вида фруктов представлены, и представители любых двух видов где-то лежат рядом. Какое наименьшее количество фруктов могло быть выложено?

\problem{2}
У сороконожки есть 40 одинаковых носков и 40 одинаковых ботинок. Она может надевать их в любом порядке с одним лишь условием: на каждую ногу надо сначала надеть носок, а только потом ботинок. Сколькими способами она может обуться?

\problem{3}
На шахматной доске стояли 16 королей, причём каждый из них бил хотя бы одного другого короля. Некоторых королей сняли с доски. Оказалось, что никакие два из оставшихся королей не бьют друг друга. Какое наименьшее количество королей могли убрать?

\problem{4}
Можно ли клетчатую доску $21 \times 21$ с вырезанными угловыми и центральной клетками разрезать на «доминошки» из двух клеток и «кресты» из пяти клеток?

\problem{5}
Для какого наименьшего $n$ на доске $10 \times 10$ можно расставить прямоугольники $1 \times 1,1 \times 2, \ldots, 1 \times n$ (возможно, повёрнутые на $90^{\circ}$ ) так, чтобы на доске не осталось места для прямоугольника $1 \times(n+1)$ ?

\problem{6}
На изначально белой доске $5 \times 5$ Петя последовательно закрашивает клетки в чёрный цвет. Внутри каждой только что закрашенной клетки он пишет число, равное количеству соседних по стороне с ней клеток чёрного цвета в этот момент. Какое наибольшее значение может принимать сумма всех выписанных чисел?

\problem{7}
У лаборанта 27 гирек. Масса каждой - целое число граммов (возможно, отрицательное). Оказалось, что если отложить любую гирю, то остальные гири можно разделить на две равные группы по 13 гирь с равными массами. Докажите, что массы всех гирь равны.

\problem{8}
На шахматную доску поставили 4 королей, не бьющих друг друга. Какое наибольшее количество королей ещё можно гарантированно поставить на доску так, чтобы все поставленные короли не били друг друга?